
\section{Scientific Qualifications}
\cvitem{Google Scholar}{\url{https://scholar.google.com/citations?user=jLujPGQAAAAJ}}

\cvitem{Publications}{I have published 19 conference papers and 8 journal papers. Additionally, two journal papers are in review, and one journal paper is in press. Since I was hired as an assistant professor at Chalmers in October 2021, I have published four journal papers on Frontiers in Computer Science, Leonardo (MIT Press), Humanities and Social Sciences Communications (Springer Nature), and Organized Sound (Cambridge Press), and 9 conference papers at conferences Sound and Music Computing Conference, New Interfaces for Musical Expression Conference, Nordic-CHI conference, ACM conference on Movement Computing, and AI Music Creativity conference.}

\cvitem{Funding}{My faculty position included a WASP-HS project (12 million SEK) as my starting package where I am the PI. Since I was hired at Chalmers, I have acquired two additional large scale projects in addition to my starting package. Along with my colleagues at Uppsala University and KTH Royal Institute of Technology, I (co-PI) have acquired VR Research Environment Grant for Social Sciences and Humanities (18 million SEK). Lastly, my colleagues from Uppsala University, KTH Royal Institute of Technology, Umeå University, Linnaeus University and I (co-PI) have acquired a WASP-HS research cluster (30 million SEK).}

\cvitem{}{Previous to my faculty position at Chalmers, I have acquired 8 grants from Canada Council for the Arts and one grant from British Columbia Arts Council in Canada as the PI; one grant from Swiss National Science Foundation and one grant from Natural Sciences and Engineering Research Council of Canada where I am the first author; one grant from Social Sciences and Humanities Research Council of Canada as the co-author.}

\section{Pedagogical Qualifications}

\cvitem{Pedagogical Training}{My pedagogical training is finalized as of June 2024 including teaching and learning in higher education (PIL 101, PIL 102, PIL 103), supervision (PIL 201) from Gothenburg University, and diversity course (CLS 930) from Chalmers, summing to a total of 22 hp ECTS credits in pedagogical training.}

\cvitem{Courses}{I have been teaching in Higher Education since 2015, and involved in 11 courses in total. At Chalmers since 2021, I have taught 1 course as the examiner (for two runs), 4 courses as the co-teacher. Prior to my faculty position at Chalmers, I taught five courses as the teaching assistant, and one course as the Sessional Instructor at Simon Fraser University in Vancouver, BC, Canada.}

\cvitem{Curriculm Development}{I have initiated a new course TRA 385 - Machine Learning and AI through Artistic Innovation at Chalmers. Previously, I have been teaching various courses covering topics in sound design, machine learning and AI, interaction design, electronics and prototyping, and art and technology. I have an extensive public teaching outreach activities outside of the university spanning twelve countries across Europe, North America, Africa, and Asia.}

\cvitem{Supervision}{I am currently the main supervisor of Kelsey Cotton (PhD student) and Xuechen Liu (Post-doctoral Fellow). Kelsey Cotton's has successfully passed her licentiate defense. I am currently hiring a new PhD student. I was one of the co-supervisors of Georgios Diapoulis who graduated in Fall 2024. Lastly, I supervised 9 masters thesis, three international summer interns (masters level), and examined 4 masters thesis.}

\section{Utilization}

\cvitem{Research Outreach}{Since 2021, I have shared my research in public events of workshops, seminars, and panels at 12 events in 7 countries in 2 continents.}

\cvitem{Artistic Outreach}{Since 2021, I have performed or exhibited my artworks at 10 public events in 7 countries in 3 continents, reaching over 60k in-person visitors/audience, with press coverage. I am the only author in these artworks, which is a common practice in audiovisual arts in live performance and exhibitions.}

\section{Academic Citizenship}
\cvitem{Internal}{At Data Science and AI (DSAI) division of CSE, I have been involved in two initiatives: 1- I have led the update of the DSAI website, 2- I joined Matti Karppa (faculty) and Emilio Jorge (doctoral student) to pursue a new computing cluster for teaching activities at DSAI. The second is now resulted in a purchase of new computing cluster, which will benefit almost all courses of DSAI. I have represented Computer Science and Engineering Department on Chalmers Day in 2022. I have organized two PhD seminars with national and international attendance, and supported those events with a WASP-HS grant. }
\cvitem{Organization}{Since October 2021, I have organized two PhD seminars in Gothenburg and Stockholm funded by WASP-HS cross-collaboration grant. I have co-organized a workshop at Nordic-CHI. I was the Arts Chair for TEI 2023 and we organized an arts exhibition at this conference. I was the industry chair of AIMC 2024.}
\cvitem{Reviewing}{Since October 2021, I have reviewed 45 conference papers, and 7 journal papers; at the venues: Frontiers in Computer Science (Springer), International Conferences on Computational Creativity, AI Music Creativity conferences, International Conferences on New Interfaces for Musical Expression, EvoMUSART, the journal Personal and Ubiquitous Computing (Springer Press), Leonardo Journal (MIT Press), Music \& Science (Sage), and Neural Computing and Applications (Springer). I had three occasions where I joined a reviewing jury of grant applications internationally: 1- I reviewed 96 grant applications for Concept to Realization Track of the Canada Council for the Arts In 2023, 2- I reviewed a grant application with the project budget of 700 000 Euros for the KU Leuven University Research Council in Belgium in 2024, 3- I reviewed a grant application at Swiss National Science Foundation.}

\cvitem{Thesis Examination}{I was the external examiner of two master thesis on music technology at University of Oslo.}